% Problem record op_8418fbac4739a7b6. This TeX file is derived from
% database/problems_json/op_8418fbac4739a7b6.json by scripts/sync-tex.mjs; edit the JSON record
% and rerun the script rather than editing this file.
\section{Minimum rank of NPT bound-entangled states}
\label{sec:op_8418fbac4739a7b6}

\paragraph{Problem.}
What is the smallest rank of a finite-dimensional bipartite state with negative
partial transpose (NPT) that is undistillable under two-way local operations and classical communication
(LOCC)? Let $\rho_{AB}$ be a density operator on
$\mathbb C^{d_A}\otimes\mathbb C^{d_B}$, with finite integers $d_A,d_B\geq2$,
$\rho_{AB}\succeq0$, and $\operatorname{Tr}\rho_{AB}=1$.
Write $\rho_{AB}^{T_B}=(\operatorname{id}_A\otimes T_B)(\rho_{AB})$ for
partial transposition on $B$. The state is NPT when
$\rho_{AB}^{T_B}\not\succeq0$.

A fixed state is undistillable precisely when, for every integer $n\geq1$
and every unit vector $\lvert\psi_n\rangle$ of Schmidt rank at most two
across the bipartition $A^n:B^n$, it satisfies
\begin{equation}
  \langle\psi_n\rvert(\rho_{AB}^{T_B})^{\otimes n}\lvert\psi_n\rangle\geq0.
  \label{eq:a7b6-all-copy}
\end{equation}
Determine the exact minimum
\begin{equation}
  r_{\min}:=\min_{\rho_{AB}}\operatorname{rank}(\rho_{AB}),
  \label{eq:a7b6-minimum-rank}
\end{equation}
where the minimization ranges over all NPT density operators satisfying
Eq.~\eqref{eq:a7b6-all-copy}, allowing arbitrary finite local dimensions.
The rank in Eq.~\eqref{eq:a7b6-minimum-rank} is the operator rank of the
joint density matrix. A complete solution must exhibit a state attaining
the proposed minimum and prove that every NPT state of smaller rank is
distillable, in every finite-dimensional bipartite system.

\subsection*{Status}

Unsolved

\subsection*{Source}

This minimum-rank formulation is derived from the low-rank obstruction and
rank-five conjecture of Chen and \DJ okovi\'c
\sourcecite{ref:a7b6-chen-djokovic}{CD16}, Theorem 5, Corollary 6, and
Conjecture 9, together with the NPT bound-entangled construction in
Faisal's Theorem 1 \sourcecite{ref:a7b6-faisal}{Fai26}.
The criterion in Eq.~\eqref{eq:a7b6-all-copy} is the finite-copy
distillability criterion of Horodecki, Horodecki, and Horodecki
\sourcecite{ref:a7b6-horodecki}{HHH98}.

\subsection*{Progress}

\begin{itemize}
  \item Chen and \DJ okovi\'c proved that every rank-four NPT state is
one-copy distillable; together with the lower-rank cases, this gives
\begin{equation}
  \operatorname{rank}(\rho_{AB})\leq4,\quad
  \rho_{AB}^{T_B}\not\succeq0
  \quad\Longrightarrow\quad
  \rho_{AB}\text{ is one-copy distillable}.
  \label{eq:a7b6-low-rank}
\end{equation}
Thus Eq.~\eqref{eq:a7b6-low-rank} implies $r_{\min}\geq5$
\sourcecite{ref:a7b6-chen-djokovic}{CD16}, Theorem 5 and Corollary 6.

  \item Chen and \DJ okovi\'c also constructed rank-five two-qutrit NPT
states $\rho_\epsilon$ in Section III. Their Lemma 8 establishes
\begin{equation}
  \forall\,n\geq1\ \exists\,\delta_n>0\ \forall\,\epsilon\in(0,\delta_n]:
  \quad\rho_\epsilon\text{ is }n\text{-undistillable}.
  \label{eq:a7b6-fixed-copy}
\end{equation}
Equation~\eqref{eq:a7b6-fixed-copy} permits the parameter to depend on
the copy number. Their Conjecture 9 asks for a fixed positive parameter
that gives undistillability for every copy number; such a state would
attain $r_{\min}=5$
\sourcecite{ref:a7b6-chen-djokovic}{CD16}.

  \item For a different family of symmetric rank-five two-qutrit states,
Lei, Song, Chen, and Liu proved one-copy undistillability on the NPT
parameter interval $[(24\sqrt2-33)/7,(33-12\sqrt6)/25)$ in Proposition 7.
Proposition 8 restricts possible two-copy negative-expectation witnesses:
a Schmidt-rank-at-most-two witness must have a component outside a
specified $17$-dimensional subspace. These results establish neither
two-copy nor all-copy undistillability
\sourcecite{ref:a7b6-lei-et-al}{LSC+26}.

  \item Faisal's Theorem 1, in arXiv preprint version 1 submitted on
25 November 2026, supplies the fixed NPT state
\begin{equation}
  \rho_* = \frac{23I-3F}{1448}
  \quad\text{on }\mathbb C^8\otimes\mathbb C^8,
  \qquad F\lvert x\rangle\lvert y\rangle=\lvert y\rangle\lvert x\rangle,
  \label{eq:a7b6-upper-example}
\end{equation}
and proves its undistillability under two-way LOCC for all copy numbers
\sourcecite{ref:a7b6-faisal}{Fai26}.
The state in Eq.~\eqref{eq:a7b6-upper-example} has eigenvalues $20/1448$
and $26/1448$ on the symmetric and antisymmetric subspaces, of dimensions
$36$ and $28$, respectively. It therefore has rank $64$ and yields
$r_{\min}\leq64$. Its partial transpose has a negative eigenvalue
$-1/1448$.
\end{itemize}

\subsection*{References}

\begin{enumerate}[label={},leftmargin=4.8em,labelsep=0.5em]
  \footnotesize
  \item[\textup{[HHH98]}]\label{ref:a7b6-horodecki}
  M. Horodecki, P. Horodecki, and R. Horodecki,
  ``Mixed-State Entanglement and Distillation: Is There a ``Bound''
  Entanglement in Nature?''
  \emph{Physical Review Letters} \textbf{80}, 5239--5242 (1998).
  \href{https://doi.org/10.1103/PhysRevLett.80.5239}{doi:10.1103/PhysRevLett.80.5239};
  \href{https://arxiv.org/abs/quant-ph/9801069}{arXiv:quant-ph/9801069}.

  \item[\textup{[CD16]}]\label{ref:a7b6-chen-djokovic}
  L. Chen and D. \v{Z}. \DJ okovi\'c,
  ``Distillability of Non-Positive-Partial-Transpose Bipartite Quantum States
  of Rank Four,'' \emph{Physical Review A} \textbf{94}, 052318 (2016).
  \href{https://doi.org/10.1103/PhysRevA.94.052318}{doi:10.1103/PhysRevA.94.052318};
  \href{https://arxiv.org/abs/1602.04416}{arXiv:1602.04416}.

  \item[\textup{[LSC+26]}]\label{ref:a7b6-lei-et-al}
  Y. Lei, Z. Song, L. Chen, and M. Liu,
  ``Entanglement Distillation of Some Rank-Five Symmetric NPT States in
  Two-Qutrit Systems,'' arXiv:2608.03710 (2026).
  \href{https://doi.org/10.48550/arXiv.2608.03710}{doi:10.48550/arXiv.2608.03710};
  \href{https://arxiv.org/abs/2608.03710}{arXiv:2608.03710}.

  \item[\textup{[Fai26]}]\label{ref:a7b6-faisal}
  I. Faisal,
  “Bound Entanglement with Negative Partial Transpose Exists,”
  arXiv:2611.34921v1 (2026), preprint, version 1, submitted 25 November 2026.
  \href{https://doi.org/10.48550/arXiv.2611.34921}{doi:10.48550/arXiv.2611.34921};
  \href{https://arxiv.org/abs/2611.34921v1}{arXiv:2611.34921v1}.
\end{enumerate}

\subsection*{Comment}

The cited lower-rank theorem and explicit construction give
\begin{equation}
  5\leq r_{\min}\leq64.
  \label{eq:a7b6-bounds}
\end{equation}
The exact value within Eq.~\eqref{eq:a7b6-bounds} remains unresolved.
In particular, the existence of a rank-five NPT bound-entangled state
would attain the lower bound; excluding rank five alone would still leave
the higher ranks to determine.

The related \href{https://qiqc-op.com/problem/op_74c4ffb5b042baf2/}{existence
of NPT bound entanglement} problem asks whether any such state exists.
The present question asks for the least possible rank. It imposes no
Werner symmetry or fixed local dimension. Proving a minimum within a
particular family or dimension, or verifying only finitely many tensor
powers, does not settle this unrestricted minimum.

\subsection*{Field}

Quantum Resource Theory

\subsection*{Topic}

Entanglement distillation; Bound entanglement; Local operations and classical communication

\subsection*{ID}

\texttt{op\_8418fbac4739a7b6}
